%-----------------------------------------------------------------------
\subsection{Instance 1: Feature Attribution Under Collinearity}
\label{sec:instance-attribution}
%-----------------------------------------------------------------------

We instantiate the abstract framework (Definition~\ref{def:explanation-system}) for
feature attribution, recovering the Attribution Impossibility of the companion paper
as a corollary of Theorem~\ref{thm:universal}.

\paragraph{The explanation system.}
Fix $P \geq 2$ features partitioned into $L$ collinearity groups.  Within each group
$\ell \in [L]$, every pair of features shares pairwise correlation $\rho \in (0,1)$;
features across groups are independent.  The explanation system
$\mathcal{S}_{\mathrm{attr}} = (\Params,\sH,Y,\mathsf{obs},\mathsf{exp},\incomp)$ is:

\begin{itemize}
  \item $\Params$: joint data-generating distributions over $\R^P$ with the collinearity
    structure above (the ground truth the data encode).

  \item $\sH$: trained models $f$ from a fixed model class (e.g.\ \GBDT{}, Lasso,
    neural networks), indexed by random training seed; each $f$ is associated with its
    data-generating distribution $\theta(f) \in \Params$ via the attribution map
    $\theta: \sH \to \Params$.

  \item $Y = (\mathrm{Fin}\,P \to \R)$: the space of feature-attribution vectors
    $\boldsymbol{\varphi} = (\varphi_1,\ldots,\varphi_P)$, where $\varphi_j(f) \geq 0$
    is the global importance of feature $j$ in model $f$ (e.g.\ mean absolute \SHAP{}
    value \citep{lundberg2017unified}, gradient-based attributions
    \citep{sundararajan2020many}, or LIME weight \citep{ribeiro2016should}).

  \item $\mathsf{obs}: \sH \to Y$ assigns to each model $f$ the attribution vector it
    ``honestly reports'': $\mathsf{obs}(f) = \boldsymbol{\varphi}(f)$.

  \item $\mathsf{exp}: \sH \to Y$ is the attribution procedure under evaluation---the
    practitioner's choice of explanation method applied to any model in $\sH$.

  \item $\incomp$: the \emph{incompatibility relation} on $Y$ captures strict reversal
    of a feature pair.  Formally, $\boldsymbol{\varphi} \incomp \boldsymbol{\varphi}'$
    if there exists a group $\ell$ and distinct features $j, k \in \ell$ such that
    \[
      \varphi_j > \varphi_k \quad \text{and} \quad \varphi'_k > \varphi'_j.
    \]
    Two attribution vectors are incompatible when they rank the same collinear pair in
    opposite orders; both orderings cannot simultaneously be correct.
\end{itemize}

\paragraph{Rashomon property.}
Collinear features induce a Rashomon set \citep{rudin2024amazing,fisher2019models,
damour2022underspecification} of near-optimal models that differ only in \emph{which}
feature within a correlated group acts as the primary predictor.
Formally, $\mathcal{S}_{\mathrm{attr}}$ has the Rashomon property
(Definition~\ref{def:rashomon}): for any group $\ell$ and any distinct features
$j, k \in \ell$, there exist models $f, f' \in \sH$ with $\theta(f) = \theta(f')$ (same
data-generating distribution) such that
\[
  \varphi_j(f) > \varphi_k(f) \qquad \text{and} \qquad
  \varphi_k(f') > \varphi_j(f').
\]
This is the \texttt{RashimonProperty} of \texttt{Trilemma.lean}: DGP symmetry within a
group guarantees every feature can serve as the first-mover in sequential optimization
(surjectivity axiom), so models with opposite collinear rankings both exist and achieve
the same training loss up to $O(\rho^2)$ \citep{damour2022underspecification,fisher2019models}.

\begin{definition}[Incompatible attribution rankings]
\label{def:attr-incompatible}
Two attribution vectors $\boldsymbol{\varphi}, \boldsymbol{\varphi}' \in Y$ are
\emph{incompatible} if there exist distinct features $j, k$ in the same collinearity
group such that $\varphi_j > \varphi_k$ and $\varphi'_k > \varphi'_j$.
\end{definition}

\paragraph{Instance 1 of the Universal Theorem.}

\begin{corollary}[Attribution Impossibility]
\label{cor:attribution}
The explanation system $\mathcal{S}_{\mathrm{attr}}$ has the Rashomon property.
Therefore, by Theorem~\ref{thm:universal}, no attribution map $\mathsf{exp}$ can
simultaneously be faithful (attributions agree with the model's own values), stable
(same attributions regardless of which equivalent model is explained), and decisive
(every collinear pair is strictly ranked).
\end{corollary}

\begin{proof}
The Rashomon property was established above.  The impossibility is then a direct
application of Theorem~\ref{thm:universal} to $\mathcal{S}_{\mathrm{attr}}$.
\end{proof}

In the notation of the companion paper (the Attribution Impossibility paper), faithfulness
corresponds to faithfulness (Definition~\ref{def:faithful}), stability to stability
(Definition~\ref{def:stable}), and decisiveness to completeness (every feature pair is
strictly ranked).  Corollary~\ref{cor:attribution} is the abstract-framework restatement
of the companion Attribution Impossibility (Theorem~1 of the companion paper), verified as
\texttt{attribution\_impossibility} in \texttt{Trilemma.lean} and as
\texttt{attribution\_impossibility\_abstract} in \texttt{AttributionInstance.lean} (zero
model-specific axiom dependencies in both cases).

\paragraph{Quantitative depth.}
The abstract impossibility is qualitative.  Architecture-specific structure supplies
quantitative bounds on \emph{how badly} the violation occurs.  For \GBDT{} with
per-group correlation $\rho$, the attribution ratio (dominant versus symmetric partner)
diverges as $1/(1-\rho^2)$ as $\rho \to 1$ (Lean-verified in \texttt{Ratio.lean}; see
also \citet{damour2022underspecification}).  For Lasso at any $\rho > 0$, one correlated
feature is zeroed out entirely, yielding an infinite ratio
(Lean-verified in \texttt{Lasso.lean}).
For random forests with $T$ trees, ensemble averaging drives the ratio to
$O(1/\sqrt{T})$---the lone converging case and the mechanism behind \DASH{}.

\paragraph{Connection to variance inflation.}
The attribution ratio $1/(1{-}\rho^2)$ coincides with the classical variance inflation factor (VIF) from regression diagnostics. This is not coincidental: both measure the amplification of estimation variance due to collinearity. The novelty is the \emph{application}: VIF measures instability of coefficient estimates, while $1/(1{-}\rho^2)$ measures instability of \emph{explanations of predictions}. These are different objects---a model's coefficients and its SHAP attributions need not agree---but the shared formula reflects the shared mechanism.

\paragraph{Resolution.}
The $G$-invariant resolution (Section~\ref{sec:resolution}) for this instance is
\DASH{} (\textbf{D}iversified \textbf{A}ggregation of \textbf{SH}AP):
averaging attribution vectors over an ensemble of
retrained models collapses the Rashomon set to a single equitable vector with
$\E[\varphi_j] = \E[\varphi_k]$ for symmetric features $j, k$, restoring faithfulness
(in expectation) and stability at the cost of decisiveness---collinear pairs are tied
rather than arbitrarily ranked.

\paragraph{Experimental note.}
Four of twenty models in the attribution experiment achieved 72--78\% held-out
accuracy; all twenty are included in the analysis. This range reflects genuine Rashomon
set variation---functionally near-equivalent models that nonetheless differ in their
attribution vectors, precisely the phenomenon the impossibility theorem characterizes.
